\section{Conclusion}

In our thesis we examined how weakening excluded-instrument relevance changes the finite-sample performance and informativeness of the preserved DML-IVQR procedure, and whether those changes differ across structural quantiles. Its contribution is a controlled Monte Carlo extension of the framework in \textcite{chen_huang_tien_2021}, rather than a new estimator or theory of inference. The design reduces the excluded signal through the parameter $\kappa$ while adding compensating noise to preserve the marginal treatment distribution and the relationship that generates endogeneity. It then compares Oracle-GMM, Full-GMM, and DML-IVQR across quantiles and sample sizes.

Point-estimation errors increase, numerical accepted sets expand, and tests generally reject false candidate effects less often. Coverage behaves differently: it can move closer to its nominal level while the test becomes less able to distinguish the truth from alternatives. The contrast is especially clear for DML-IVQR at $n=1000$, $\tau=0.50$, and $\kappa=0.10$. Coverage is 0.942, yet the median accepted-set measure is 3.900 within a domain of width four, and all 41 grid points are accepted in 0.428 of replications. These are distinct simulation summaries, and their underlying acceptance events need not coincide replication by replication. Together, however, they show why frequent retention of the truth does not by itself establish informative inference. Calibration concerns rejection at the truth; informativeness concerns the ability to exclude false values.

Oracle-GMM knows the designated relevant controls, but its accuracy, power, and numerical localisation still deteriorate as relevance weakens. High-dimensional nuisance estimation therefore cannot be the sole explanation for the common decline in performance. DML-IVQR often improves point accuracy relative to Full-GMM, showing that regularised nuisance estimation can be useful in this high-dimensional design, but the advantage is neither universal across design cells nor transferable automatically to coverage or raw power. Oracle-GMM is likewise an informative benchmark rather than a guaranteed finite-sample winner. Moreover, the three procedures differ in their complete implementations, including nuisance dimension and residualisation as well as regularisation. Their comparison therefore does not isolate the causal effect of regularisation, and differences in unadjusted power cannot be read as pure efficiency rankings when empirical sizes also differ.

Sample size and quantile comparisons reveal a similar separation among performance measures. Raising $n$ from 500 to 1000 lowers observed RMSE in every matched estimator--quantile--$\kappa$ comparison, but it does not ensure improved calibration. Several DML-IVQR coverage rates deteriorate substantially in the strong-relevance design even as point accuracy improves. Across quantiles, the final results have a clear accuracy ordering: at fixed estimator, sample size, and $\kappa$, median MAE and RMSE are lower than at either $\tau=0.10$ or $\tau=0.90$ throughout. Coverage has no corresponding universal centre-versus-tail ordering. The evidence therefore supports finite-sample statements about these two sample sizes; it neither identifies an asymptotic rate nor shows that lower RMSE reflects lower estimator variance alone.

Neyman orthogonality is intended to reduce first-order sensitivity to sufficiently small nuisance-estimation errors \parencite{chernozhukov_et_al_2018}. Candidate-specific test construction is intended to accommodate weak identification under its assumptions without relying on a precisely estimated structural coefficient \parencite{chernozhukov_hansen_2008}. Neither property creates identifying variation, guarantees accurate finite-sample critical values, or ensures that false candidate effects can be excluded. The simulations illustrate the practical importance of evaluating nuisance robustness, finite-sample calibration, power, and accepted-set informativeness separately. They do not establish or refute the general asymptotic theory of DML-IVQR; they evaluate the executable procedure preserved for this experiment.

That scope matters because the relationship between the executable procedure and the ideal construction is qualified in several ways. For this DGP, the population density weight at the true parameter encounters the non-finite regularity problem derived in Section~3.2. The implemented procedure instead uses a fitted-normal density proxy and implementation-specific residualisation rules, as described in Section~3.3. Appendix~\ref{app:reproducibility} records that the preserved \texttt{dnorm} call supplies a residual variance in the argument position interpreted by R as a standard deviation. The DML implementation estimates nuisance functions and scores on the same sample rather than using cross-fitting, and fresh pivotal draws introduce algorithmic randomness across candidate evaluations. All three estimators use an asymptotic chi-square cutoff even though their empirical sizes differ. These features may matter for finite-sample behaviour, but the experiment does not vary them one at a time. It therefore cannot attribute any particular distortion to one implementation choice or establish that changing that choice would restore calibration.

Point estimates are first minimisers on a 41-point grid and consequently reflect its resolution, endpoints, and tie rule. The main accepted-set measure describes only the maintained domain $A_0=[-1,3]$. The domain-expansion exercise in Appendix~\ref{app:numerical-domain-diagnostics} shows that acceptance can continue over wider investigated domains when relevance is weak, but continued acceptance on a finite sequence of grids is not proof of global unboundedness or non-identification. More broadly, the results come from one DGP family, selected relevance settings, two sample sizes, and five quantiles. They document robust patterns within that design but cannot determine how every alternative DGP, implementation, or larger sample would behave.

These limitations point to focused extensions. First, controlled variants could alter the density proxy, residualisation, sample splitting, or treatment of algorithmic randomness one feature at a time while retaining the present implementation as an explicit baseline and applying common evaluation criteria. Such comparisons would assess sensitivity and possible mechanisms without presuming that any single modification must repair the observed calibration. Second, a complementary DGP satisfying the relevant density-weight regularity requirements could help distinguish properties of the present simulation environment from the behaviour of the theoretical procedure under its maintained assumptions. Third, additional sample sizes and relevance settings could be combined with systematic grid and domain checks. Direct-at-truth coverage should remain distinct from numerical inversion in that work: more replications improve Monte Carlo precision but do not remove systematic miscalibration, and no finite grid can establish global boundedness.
